% !PRINTED 79
\weberpage{79}

Wenn die Primzahl $q$ nur einfach in $m$ aufgeht und $\mathfrak{A}$ durch $Q$ theilbar ist, so ist zwar $\mathfrak{A}$ nicht primär; es gehört aber trotzdem die Periode $\eta$ zu der Gruppe $\mathfrak{A}$.

Es kann aber in diesem Falle $\eta$ durch Einheitswurzeln von niedrigerem Grade dargestellt werden. In folgender Weise lässt sich diese Darstellung finden. Alle Zahlen von $Q$ sind von der Form $1 + hm'$, worin $h$ die Reihe der Zahlen $0, 1, \dots q - 1$ durchläuft, mit Ausnahme des einzigen Werthes, für den

\weq{18}{1 + hm' = kq}

durch $q$ theilbar wird. Setzt man also unter der Voraussetzung, dass $\mathfrak{A}$ durch $Q$ theilbar ist,

\[
\mathfrak{A} = Q + Qa_1 + Qa_2 \dots,
\]

so folgt, dass jede Zahl $a$ in $\mathfrak{A}$ von der Form ist:

\weq{19}{a = a'(1 + hm'),}

worin $a'$ die Reihe der Zahlen $1, a_1, a_2 \dots$ durchläuft. Daraus ist noch zu schliessen, dass die Reste der Zahlen $1, a', a'' \dots$ nach dem Modul $m'$ eine Gruppe bilden. Nimmt man nun die Summe über alle Werthe $h = 0, 1, \dots q - 1$ und nimmt dann den einen durch (18) bestimmten Werth wieder weg, so folgt:

\[
\eta = \sum^a r^a = \sum^{a'} \sum^h r^{a'(1 + hm')} - \sum^{a'} r^{a'kq},
\]

oder, da $\sum^h r^{a'm'h} = 0$ ist:

\[
\eta = -\sum^{a'} r^{a'kq}.
\]

Nun ist $r^q$ eine $m'^{\text{te}}$ Einheitswurzel, und $\eta$ also, vom Vorzeichen abgesehen, gleich einer aus $m'^{\text{ten}}$ Einheitswurzeln gebildeten Periode\footnote{Hier ist die Abhandlung von \gesperrt{Fuchs} zu vergleichen: Ueber die Perioden, welche aus den Wurzeln der Gleichung $\omega^n = 1$ gebildet sind, wenn $n$ eine zusammengesetzte Zahl ist. Crelle's Journal, Bd. 61 (1863).}.

\begin{center}
\S.\,21.

\gesperrt{Kreistheilungskörper mit gegebener Gruppe}.
\end{center}

Im Vorhergehenden hat sich also ergeben, dass jede Kreistheilungsperiode

\[
\eta = \sum^a r^a,
\]

in der $r$ eine primitive $m^{\text{te}}$ Einheitswurzel ist und $a$ die Zahlen einer Gruppe $\mathfrak{A}$ durchläuft, wenn $\mathfrak{A}$ ein primärer Theiler von $\mathfrak{N}$
%CONT>>
